\section{Sliding Window Advertisement (Maximum)}
\label{sec:cses-sw-advertisement}

\problemstatement{Given an array of $n$ integers and a window size $k$,
output the maximum of every contiguous window. (This ``advertisement''
variant is the maximum counterpart of the sliding window minimum.)}

\begin{patternbox}
Maximum, like minimum, is not invertible, so the same \textbf{monotonic
deque} applies -- kept in \emph{decreasing} order this time, so the front is
always the current window's maximum. Each index is pushed and popped once,
giving $O(n)$.
\end{patternbox}

\subsection*{A decreasing deque}

Maintain a deque of indices whose values are strictly decreasing from front
to back. When a new element arrives, pop from the back every index whose
value is $\le$ the new value -- those are dominated and can never be the
maximum again while the newer, larger element is present. Push the new
index, drop the front if it has left the window, and read the front as the
window maximum.

\begin{ideabox}[Mirror of the Minimum Deque]
The only change from sliding window minimum is the comparison direction:
pop while the back is $\le$ the incoming value (instead of $\ge$), keeping
the deque decreasing so its front is the maximum. One structure, two mirror
uses.
\end{ideabox}

\subsection*{Algorithm}

\begin{enumerate}
  \item For each index $i$: pop back while its value $\le a[i]$; push $i$.
  \item Pop front while it is $\le i-k$ (out of window).
  \item Once $i\ge k-1$, output $a[\text{front}]$.
\end{enumerate}

\subsection*{Dry run}

\dryrun{$a=[3,1,4,2]$, $k=2$.}

\begin{itemize}
  \item $[3,1]$: deque front $3$; max $3$.
  \item $[1,4]$: $4$ pops $1$ (and $3$ leaves window); max $4$.
  \item $[4,2]$: front $4$; max $4$. Output $3,4,4$.
\end{itemize}

\subsection*{C++ implementation}

\begin{lstlisting}
#include <bits/stdc++.h>
using namespace std;

int main() {
    int n, k;
    cin >> n >> k;
    vector<int> a(n);
    for (int& x : a) cin >> x;

    deque<int> dq;                            // indices, values decreasing
    for (int i = 0; i < n; i++) {
        while (!dq.empty() && a[dq.back()] <= a[i]) dq.pop_back();
        dq.push_back(i);
        if (dq.front() <= i - k) dq.pop_front();
        if (i >= k - 1) cout << a[dq.front()] << " ";
    }
    cout << "\n";
    return 0;
}
\end{lstlisting}

\begin{complexitybox}
\textbf{Time Complexity.} $O(n)$ -- each index enters and leaves the deque
once. \textbf{Space Complexity.} $O(k)$.
\end{complexitybox}

\begin{takeawaybox}
Sliding window maximum is the minimum deque mirrored: pop while the back is
no larger than the incoming value. Master the monotonic deque once and both
extremes -- and many ``nearest greater/smaller'' problems -- follow from it.
\end{takeawaybox}
